\vspace*{6cm}
\begin{center}
{\Large \gesperrt{VORWORT.}}

\vspace{1em}
\rule{2cm}{0.4pt}
\end{center}
\vspace{2em}

Es ist mir vergönnt, den Plan einer Weiterführung meines Lehrbuches der Algebra, den ich vor zwölf Jahren in der Vorrede zur ersten Auflage des zweiten Bandes angekündigt habe, nach mannigfaltigen Abhaltungen noch auszuführen. Durch das Entgegenkommen der Verlagsfirma erscheint dieser dritte Band der Algebra zugleich als zweite Auflage der im Jahre 1891 zum erstenmal gedruckten „Elliptischen Funktionen und algebraischen Zahlen“.

Er beschäftigt sich hauptsächlich mit dem weiteren Ausbau der mannigfaltigen Anwendungen der Algebra und besonders der Theorie der quadratischen Körper auf die aus den elliptischen Funktionen hervorgegangenen Probleme, die uns das erste über die Kreisteilung hinausgehende Beispiel von algebraischen Zahlen liefern, deren Gesetze einigermaßen bekannt sind. Als Grundlage dazu dient eine eingehendere Behandlung der quadratischen Körper mit negativer Diskriminante. Freilich ist auch hier nicht alles erreicht, was ich mir als letztes Ziel gesteckt hatte. So mußte die Ausführung der Theorie der relativ zyklischen Körper noch zurückgestellt werden --- hoffentlich nur einstweilen.

Dagegen habe ich, einem mehrfach an mich herangetretenen Wunsche entsprechend, einen Abriß der Theorie der algebraischen Funktionen auf arithmetischer Grundlage beigefügt, der
